\section{Conclusion and Future Work}
\label{sec:conclusion}

In this paper, we introduced \textbf{Resource-Budgeted Top-$M$ Expert Serving (RB-TopM)}, a training-free, hardware-aware framework for dynamic activation blending on the edge. Guided by the core philosophy of our Pragmatist persona, we designed RB-TopM to solve the critical serving bottleneck of modern ensembling SOTA: uncontrolled compute scaling and memory bandwidth strain on low-power devices. By introducing real-time closed-loop control functions governed by a resource availability parameter $C_{\text{budget}} \in [0, 1]$, RB-TopM dynamically scales the ensembling capacity and prunes marginal adapter branches in microseconds on-the-fly. 

We validated RB-TopM on a 14-layer Analytical Coordinate Sandbox simulating heterogeneous visual streams over 10 independent random seeds. RB-TopM successfully establishes a smooth, monotonic accuracy-latency trade-off frontier, achieving up to 78.4\% in expert FLOP savings with virtually zero degradation in task ensembling accuracy compared to un-gated baselines. Additionally, we showed that an early Coordinate GMM safety shield successfully rejects 38.04\% of high-noise OOD queries, preventing corrupted inputs from triggering downstream specialized adapter execution.

Having successfully demonstrated physical bare-metal deployment and high-resolution profiling of RB-TopM on the STM32 microcontroller using a Joulescope JS110 analyzer, our future work will focus on expanding this physical validation to high-performance edge accelerators (such as Raspberry Pi NPU boards or Nvidia Jetson Orin modules) running multi-gigabyte models under actual real-time concurrency. 

To achieve even higher overall full-system serving latency speedups, we are highly interested in combining RB-TopM with base backbone optimizations. While our method reduces expert-only execution by up to 76.2\%, the pre-trained base model backbone must execute in its entirety (consuming 4.85 ms in our MobileNetV3-Large TVM-compiled simulation pilot). Integrating base backbone early-exiting (e.g., BranchyNet, which exits early at intermediate layers under high routing confidence) or structural base weight pruning/quantization could shrink this base backbone latency, multiplying full-system speedups. Additionally, to handle non-stationary physical environments where representation drift occurs gradually (e.g., due to diurnal lighting shifts or camera lens degradation), we plan to implement an online centroid adaptation mechanism. By performing microsecond-scale Exponential Moving Average (EMA) updates on task centroids using highly confident, in-distribution test-time samples ($\mu_k^{(3)} \gets (1 - \lambda) \mu_k^{(3)} + \lambda h_b^{(3)}$), the router can adapt continuously to representation drift and maintain long-term serving robustness in the wild.

To address the GMM safety shield's calibration generalization trade-off (where regularizing calibration data via 5-fold cross-validation corrects test-set false positives but reduces OOD detection sensitivity on high-noise queries), we are eager to investigate highly compact, integer-friendly \textbf{Normalizing Flows} (such as lightweight RealNVP or Masked Autoregressive Flows with lookup-table approximations) tailored for edge platforms. Normalizing flows can learn complex, tight, and non-linear task similarity boundaries without requiring extensive mixture parameters, successfully preserving outstanding OOD rejection rates under regularized settings. Furthermore, to bridge the ensembling accuracy gap on highly overlapping visual domains like CIFAR-10 (where representation bleed and early-stage routing misclassifications lead to a drop below the Expert Oracle), we propose transitioning from linear cosine similarity projections to small, non-linear bottleneck projection layers. A lightweight, non-linear routing head (such as a two-layer Multi-Layer Perceptron with ReLU activations) can significantly enhance the angular and spatial separation of complex, high-dimensional representation manifolds, preventing activation-space bleeding and bringing ensembling performance closer to the Expert Oracle ceiling.

Furthermore, we are highly interested in implementing OS-level drivers that hook the $C_{\text{budget}}$ parameter directly into Linux kernel system events (e.g., \texttt{sysfs} power supply states, CPU thermal throttling interrupts, or Android battery saver broadcasts). Our formal systems-level analysis (Appendix~\ref{app:closed_loop_mapping}) demonstrates that evaluating our closed-loop controller at 100 Hz consumes less than $4.79 \times 10^{-5}\%$ of CPU capacity on low-power microcontrollers and only $0.012\%$ of a single CPU core on Linux edge boards. A 100 Hz polling frequency is highly sufficient because thermal throttling and battery depletion are physical macro-phenomena that unfold over seconds or minutes, and queue depth averages are highly stable over 10 ms. Under extremely volatile workloads or bursty traffic patterns, practitioners can scale the polling frequency up to 1 kHz (1 ms resolution) to prevent serving queue latency spikes, while the CPU overhead remains entirely negligible ($<0.001\%$ of CPU capacity). This confirms that the practical systems-level deployment overhead of the closed-loop driver is entirely negligible, enabling sustainable, real-time hardware-guided routing in volatile, production edge environments.
